% Panduan Operasional ISHAS — sumber LaTeX.
% Kompilasi: latexmk -pdf panduan-operasional.tex
\documentclass[11pt,a4paper]{article}

\usepackage[T1]{fontenc}
\usepackage[utf8]{inputenc}
\usepackage{underscore}
\usepackage[english,indonesian]{babel}
\usepackage{lmodern}
\usepackage{microtype}
\usepackage[a4paper,margin=2.4cm,footskip=1cm]{geometry}
\usepackage[table]{xcolor}
\usepackage{booktabs}
\usepackage{array}
\usepackage{enumitem}
\usepackage{amssymb}
\usepackage{listings}
\usepackage{titlesec}
\usepackage{fancyhdr}
\usepackage[hidelinks]{hyperref}

\definecolor{ishas}{HTML}{9F1239}
\definecolor{ishaslight}{HTML}{FCE7F3}
\definecolor{codebg}{HTML}{F6F7F9}
\definecolor{codeframe}{HTML}{D1D5DB}

\titleformat{\section}{\color{ishas}\Large\bfseries}{\thesection}{0.6em}{}
\titleformat{\subsection}{\color{ishas}\large\bfseries}{\thesubsection}{0.6em}{}
\titleformat{\subsubsection}{\bfseries}{\thesubsubsection}{0.6em}{}

\pagestyle{fancy}
\fancyhf{}
\fancyhead[L]{\small\textcolor{ishas}{Panduan Operasional ISHAS}}
\fancyhead[R]{\small\thepage}
\renewcommand{\headrulewidth}{0.4pt}

\lstset{
  basicstyle=\ttfamily\small,
  breaklines=true,
  breakatwhitespace=false,
  columns=fullflexible,
  frame=single,
  rulecolor=\color{codeframe},
  backgroundcolor=\color{codebg},
  keepspaces=true,
  showstringspaces=false,
  xleftmargin=0.6em,
  xrightmargin=0.2em,
  aboveskip=0.6em,
  belowskip=0.6em
}

\setlist{topsep=2pt,itemsep=2pt,parsep=0pt}
\newcommand{\code}[1]{\texttt{#1}}

\title{\color{ishas}Panduan Operasional ISHAS}
\author{Integrated Safety and Health Assessment System For Boarding School}
\date{\today}

\begin{document}

\begin{titlepage}
  \centering
  \vspace*{3cm}
  {\color{ishas}\Huge\bfseries Panduan Operasional ISHAS\par}
  \vspace{0.6cm}
  {\large Integrated Safety and Health Assessment System For Boarding School\par}
  \vspace{0.4cm}
  {\large Pemasangan, data, cadangan, pembaruan, dan serah terima\par}
  \vspace{2cm}
  {\large Versi dokumen: \today\par}
  \vfill
  \begin{minipage}{0.86\textwidth}
    \colorbox{ishaslight}{%
      \parbox{\dimexpr\linewidth-2\fboxsep}{%
        \textbf{Untuk siapa:} operator/teknisi yang memasang dan merawat aplikasi.
        Bukan untuk pengguna harian --- lihat \emph{Panduan Pengguna}.
        \textbf{Status:} prototipe; data contoh, rumus/ambang ilmiah final belum
        ditetapkan.%
      }%
    }
  \end{minipage}
  \vspace{1.2cm}
\end{titlepage}

\tableofcontents
\newpage

\section{Prasyarat}
\begin{itemize}
  \item Mesin dengan \textbf{Docker + Docker Compose}.
  \item Salinan repository ISHAS.
  \item (Opsional, tanpa Docker) \textbf{Bun} dan \textbf{MySQL 8} langsung di mesin.
  \item File \code{.env} di root (salin dari contoh lalu isi).
\end{itemize}

\begin{lstlisting}
cp .env.example .env
\end{lstlisting}

\noindent Nilai penting di \code{.env}: \code{DB_NAME}, \code{DB_USER},
\code{DB_PASSWORD}, \code{DB_ROOT_PASSWORD}, \code{WEB_PORT}, \code{API_PORT},
\code{SEED_DEFAULT_PASSWORD}, \code{VITE_USE_BACKEND}, dan \code{COOKIE_SECURE}
(isi \code{true} bila sudah di balik HTTPS). Periksa hasil pembacaan:
\code{docker compose config}.

\section{Menjalankan aplikasi}

\subsection{Mode ngoding (dev) --- satu perintah}
\begin{lstlisting}
docker compose --profile dev up --build
\end{lstlisting}
Menyalakan \textbf{database + backend + frontend} sekaligus. Buka
\code{http://localhost:3003} (login memakai kartu akun contoh).

\subsection{Mode rilis (prod)}
\begin{lstlisting}
docker compose --profile prod up --build
\end{lstlisting}
Frontend statis (nginx) + backend rilis; halaman login memakai \textbf{email +
kata sandi}, tombol akun demo dimatikan.

\subsection{Tanpa Docker (khusus backend)}
\begin{lstlisting}
cd apps/api
bun install
bun run migrate
bun run seed --mode=demo     # atau --mode=empty
bun run dev                  # http://localhost:3004
\end{lstlisting}

\section{Menyiapkan / mengulang data}

Urutan yang disarankan: migrasi schema, lalu isi data.

\begin{lstlisting}
# via Docker (container rilis)
docker compose --profile prod run --rm api-prod bun run migrate --fresh
docker compose --profile prod run --rm api-prod bun run seed --mode=demo   # atau --mode=empty

# via host
cd apps/api
bun run migrate              # tambah/tambahkan migrasi
bun run migrate --fresh      # bangun ulang schema dari nol (menghapus data)
bun run seed --mode=demo     # data contoh lengkap untuk demo
bun run seed --mode=empty    # data inti/kosong untuk mulai bersih
\end{lstlisting}

\begin{itemize}
  \item \code{--mode=demo}: mengisi data contoh (pesantren, laporan, instrumen, dsb).
  \item \code{--mode=empty}: hanya data inti (akun dasar + bank instrumen) untuk
        mulai bersih.
  \item \textbf{Reset data demo} juga tersedia dari UI: \code{/admin/pengaturan}
        $\rightarrow$ \emph{Reset data demo} (memuat ulang data contoh).
\end{itemize}

\section{Akun \& kata sandi}
\begin{itemize}
  \item Kata sandi awal akun seed = \code{SEED_DEFAULT_PASSWORD} di \code{.env}.
  \item Pengguna mengganti sandi sendiri lewat menu akun $\rightarrow$
        \emph{Ganti kata sandi}.
  \item Super Admin dapat \textbf{reset sandi} akun lain di
        \code{/admin/pengaturan} pengguna.
  \item Jangan menaruh sandi/token produksi di dalam kode atau di browser.
\end{itemize}

\section{Pemeriksaan kesehatan}
\begin{itemize}
  \item Backend: \code{GET http://localhost:3004/health} (juga tersedia di
        \code{/api/v1/health}).
  \item Lewat web (di balik nginx): \code{GET http://localhost:3003/api/v1/health}.
  \item Uji cepat setelah pasang: halaman \code{/} terbuka, \code{/login} tampil,
        dan login akun contoh berhasil.
\end{itemize}

\section{Pencadangan (backup) \& pemulihan}

Cadangkan \textbf{dua hal}: basis data dan berkas unggahan.

\subsection{Basis data}
\begin{lstlisting}
# dump ke file di host
docker compose exec db sh -c \
  'mysqldump -uroot -p"$MYSQL_ROOT_PASSWORD" "$MYSQL_DATABASE"' \
  > ishas-db-$(date +%F).sql
\end{lstlisting}
Pemulihan:
\begin{lstlisting}
cat ishas-db-YYYY-MM-DD.sql | docker compose exec -T db sh -c \
  'mysql -uroot -p"$MYSQL_ROOT_PASSWORD" "$MYSQL_DATABASE"'
\end{lstlisting}

\subsection{Berkas unggahan (volume \code{api-storage})}
Nama volume mengikuti nama proyek Compose; temukan dulu:
\begin{lstlisting}
docker volume ls | grep api-storage
\end{lstlisting}
Cadangkan:
\begin{lstlisting}
docker run --rm \
  -v <nama_volume_api-storage>:/data \
  -v "$PWD:/backup" \
  alpine tar czf /backup/ishas-storage-$(date +%F).tar.gz -C /data .
\end{lstlisting}
Pemulihan:
\begin{lstlisting}
docker run --rm \
  -v <nama_volume_api-storage>:/data \
  -v "$PWD:/backup" \
  alpine sh -c 'tar xzf /backup/ishas-storage-YYYY-MM-DD.tar.gz -C /data'
\end{lstlisting}
Jadwalkan backup berkala (cron / \code{systemd timer}) untuk DB dan volume storage.

\section{Pembaruan (upgrade) \& rollback}
Pembaruan:
\begin{lstlisting}
git pull
docker compose --profile prod build
docker compose --profile prod run --rm api-prod bun run migrate
docker compose --profile prod up -d
\end{lstlisting}

\noindent Rencana rollback (uji asap pasca-pasang --- D4):
\begin{enumerate}
  \item Hentikan layanan: \code{docker compose --profile prod down}.
  \item Pulihkan basis data + volume dari backup terakhir (\S6).
  \item Kembalikan kode ke versi/tag sebelumnya (\code{git checkout <tag>}), lalu
        \code{docker compose --profile prod up --build}.
  \item Uji asap: \code{/health}, \code{/login}, dan alur utama (lihat \S5 dan
        \emph{Panduan Pengguna}). Catat hasilnya.
\end{enumerate}

\section{Pemeliharaan rutin}
\begin{itemize}
  \item Periksa kesehatan API dan ruang disk.
  \item Bersihkan berkas yatim (staging $>$24 jam) bila perlu:
        \code{cd apps/api \&\& bun run sweep}.
  \item Perbarui image dasar (\code{docker compose build --pull}) secara berkala.
  \item Pantau log: \code{docker compose logs -f api-prod web-prod}.
\end{itemize}

\section{Serah terima \& pelatihan (checklist)}

Dipakai saat menyerahkan aplikasi ke pemilik:
\begin{itemize}
  \item[$\square$] Berikan akses: alamat aplikasi, akun contoh, dan \emph{Panduan
        Pengguna}.
  \item[$\square$] Serahkan dokumen: \code{README.md} (root), folder \code{docs/},
        \code{.env.example}.
  \item[$\square$] Latih tiap peran: publik/pelapor, Pesantren, Validator, Super
        Admin.
  \item[$\square$] Jelaskan cara reset data demo dan mengganti kata sandi.
  \item[$\square$] Buka data pada server, lakukan uji asap, dan catat hasil (D4).
  \item[$\square$] Tandatangani berita acara serah terima (nama, peran, tanggal).
\end{itemize}

\vfill
\noindent Dukungan dokumen: \code{DEPLOYMENT\_CLOUDFLARE\_TUNNEL.md},
\code{BACKEND\_MIGRATION.md}, \code{BACKEND\_STORAGE.md}, \code{README.md} (root),
\code{apps/api/README.md}.

\end{document}
